%==============================================================================
\section{Introduction}
\label{sec:intro}

% --- maps to 01_Project_Overview.md §1, §2 ---

\subsection{Motivation and point of departure}
Integrated circuits degrade. Bias-temperature instability (BTI/NBTI), hot-carrier
injection and related mechanisms slow critical paths over a device's lifetime,
and SLM aims to monitor that degradation in-field to enable adaptive voltage
scaling and remaining-useful-life (RUL) prediction. This project builds on three
correlated prior efforts: an auto-tuning on-chip slack sensor
\cite{nogueira_sensor}, a thermal-control fidelity study comparing bang-bang and
PID burn-in ovens on the same device and sensor \cite{alencar_fidelity}, and a
low-cost automated burn-in platform \cite{alencarfilho_tcc}. Across all three, a
single fact recurs: the slack a sensor reports is a \emph{mixture} of a
permanent aging term and a reversible PVT term,
\begin{equation}
  \signalmodel ,
  \label{eq:model}
\end{equation}
where $\eq$ is quantization noise bounded by the sensor LSB
(\valLSBps~ps per count).

\subsection{The gap}
The prior work measures the reversible contamination with the linear $R^2$ of
slack regressed on $\Tdie$ and $\Vcore$ (e.g.\ \valRsqBangBang{} under bang-bang
vs.\ \valRsqPid{} under PID control). This has two limitations central to this
physical domain: it captures \emph{only linear} dependence, although BTI is
super-linear in $\Vcore$ and the Arrhenius acceleration is exponential in
inverse temperature; and it yields \emph{no transferable unit}---``variance
explained'' does not answer how many aging states the sensor resolves, how early
a trend is detectable, or how much aging information a given bench preserves.
Information theory answers all three in one currency: \textbf{bits}.

\subsection{Contribution}
We add an analytical measurement layer---no new silicon, sensor, or oven---that
re-expresses noise, correction and recoverability in the language of entropy,
mutual information and channel capacity. Concretely:
\begin{enumerate}
  \item $\MI{\slk}{\Tdie,\Vcore}$ replaces $R^2$ as a non-linear contamination
        measure (\cref{sec:results});
  \item an entropy account of PVT correction quantifies, in bits, the
        uncertainty each auxiliary channel removes;
  \item a capacity estimate and a BSC alarm model convert measured noise into
        resolvable aging states and alarm reliability;
  \item a Cram\'er--Rao/Bayesian detectability bound states the earliest
        detectable aging trend given the noise floor, complemented by a
        Wiener-process/Kalman-filter treatment that turns the same problem
        into a continuously-updated, denoised aging trajectory and a
        closed-form Inverse-Gaussian RUL \emph{distribution}
        (\cref{sec:res-wiener}).
\end{enumerate}
We additionally exercise source coding, Kolmogorov/MDL model selection and
ICA-based blind separation on the same logs, so that the project spans the full
syllabus of the course (\cref{sec:methodology} and the topic mapping).

\subsection{Scope and honesty}
The primary dataset is a single \valDurationHours~h PID campaign on one device;
it is long enough to expose a measurable aging drift but cannot, alone,
cross-validate a permanent-aging trajectory across devices. RUL results are thus
framed as a \emph{methodology and a detectability bound}. Mutual-information
estimation on strongly autocorrelated series is bias-prone; we report everything
on an effective sample size $\neff$ with block-bootstrap intervals.
